\section{The decision rule}
\label{sec:rule}
\label{sec:method}

The decision rule divides the work: a model reads the file and gives each fact a value, and code applies the workflow to those values. This section defines the rule and states four results with their scope; the proofs are in Appendix~\ref{app:proofs}.

We carry the mould claim of Section~\ref{sec:intro} throughout. One file says: ``Our daughter has been coughing for weeks; the paediatrician suspects a link with the mould. I do not have a medical report yet.'' Another says: ``\ldots we do not know whether it has anything to do with the mould and have not seen a doctor.''

\subsection{Workflows and facts}
\label{sec:workflow}

A workflow states an organisation's rules for one kind of claim. It names a finite set $\mathcal{C}$ of facts, the conditions its rules depend on, such as $\mathit{health}$: the mould has affected someone's health. A fact is true ($\T$), false ($\F$) or unresolved ($\U$).

The workflow arranges steps in an acyclic graph $G=(N,E)$ with one entry step. Each step $n$ has a condition $c_n$, and each transition $(u,n)\in E$ has a condition $c_{un}$. Each document type $d$ is needed at a set of steps $L(d)\subseteq N$ under its own condition $c_d$. Conditions are formulas over facts built with \emph{and}, \emph{or} and \emph{not}. We call such a workflow a \emph{flow workflow}: its transitions only enable or branch to a step, with no relation that blocks, suppresses or waits for another step.

The output, a checklist, marks each document type \request, \hold, \skipd or \onfile, and lists as questions the facts on which held documents wait. Whether a document is on file is decided by two model judgments (Section~\ref{sec:pipeline}); the propositions take it as given.

The product also has facts that take one of a list of values, a date or a text; for exact deferral such an open fact takes each listed value, or each value the conditions name and one for any other. In the tested workflows only the fact that selects the workflow is not two-valued. Completions treat facts as independent, so for related facts exact deferral can hold a document that every completion consistent with the relation needs.

\subsection{Readings, admission and presumptions}
\label{sec:admission}

A reading is one model pass over the file that answers every fact true or false with a quotation, open (the file says it is not known, quoting the words) or not stated. \casepath makes three readings by default, each a separate request. A quotation is \emph{found} when, after soft hyphens are removed, words split by a hyphen at a line break are joined and white space is collapsed, it occurs letter for letter, with letter case kept, on one page of a file; failing that, a second match also equates the typographic variants of quotation marks, apostrophes, dashes, ligatures and spaces. Nothing else is folded, and a quotation shorter than eight characters is never found. A value whose quotation is not found, even after one request to quote again, counts as no value.

Table~\ref{tab:admission} gives the admission rule of the tested commit for the facts that decide the checklist or the routing: those in a transition, in a document's condition or in the workflow's selection condition. Code, not the model, applies it. Other facts are recorded from the first reading alone and do not enter the evaluation.

\begin{table}[t]
\caption{Admission of a fact that decides the checklist or the routing, at commit 6a72159c. ``Value'' means a true or false answer, or a value, with its quotation found. A fact left unresolved holds the documents whose need it can change.}
\label{tab:admission}
\footnotesize
\begin{tabularx}{\linewidth}{@{}>{\raggedright\arraybackslash}X>{\raggedright\arraybackslash}p{0.26\linewidth}>{\raggedright\arraybackslash}p{0.26\linewidth}@{}}
\toprule
What the readings give & No presumption & Presumption of value $v$ \\
\midrule
Every reading gives the same value $x$ & $x$ is admitted & $x$ is admitted \\
Readings give different values & unresolved; marked for a person & unresolved; marked for a person \\
Some readings give $x$, the others no value & unresolved & $v$ if $x=v$; otherwise unresolved \\
No reading gives a value; one or more answer open with words found & unresolved & unresolved \\
No reading gives a value; none answers open with words found & unresolved & $v$ \\
\bottomrule
\end{tabularx}
\end{table}

For the first file every reading gives $\mathit{health}=\T$ and quotes ``Our daughter has been coughing for weeks; the paediatrician suspects a link with the mould'', which is found, so $\T$ is admitted. For the second every reading answers open, quoting ``we do not know whether it has anything to do with the mould''.

A file that never mentions a fact does not express doubt about it. A \emph{presumption} states, with a clause from the organisation's sources, what an unstated fact is taken to be. The checklist records a presumed value as a presumption. The tenancy workflow presumes that the mould has affected no one's health when the file never mentions health. In the second file the readings answer open, so $\mathit{health}$ stays $\U$.

Admitted values and presumptions form an assignment $\val:\mathcal{C}\to\Vals$. A \emph{completion} of $\val$ is an assignment $\val':\mathcal{C}\to\{\T,\F\}$ that agrees with $\val$ wherever $\val(f)\neq\U$.

\subsection{Evaluation}
\label{sec:kleene}

Code evaluates conditions in strong Kleene logic \cite{kleene1952introduction}. With the order $\F<\U<\T$, \emph{and} takes the minimum, \emph{or} the maximum, and \emph{not} swaps $\T$ and $\F$. A step is active when it is the entry step or is reached from an active step through a transition that holds, and its own condition holds:
\begin{equation}
\act(n) \;=\; \Bigl[\,\mathrm{entry}(n)\;\kor\;\bigvee_{(u,n)\in E}\bigl(\act(u)\kand\val(c_{un})\bigr)\Bigr]\;\kand\;\val(c_n).
\label{eq:activation}
\end{equation}
A document is needed when a step that needs it is active and its own condition holds:
\begin{equation}
\req(d) \;=\; \Bigl[\,\bigvee_{n\in L(d)}\act(n)\Bigr]\;\kand\;\val(c_d).
\label{eq:requirement}
\end{equation}
A document with $\req(d)=\F$ is \skipd; one with $\req(d)=\U$ is held; one with $\req(d)=\T$ is \onfile if the file contains it and requested otherwise. We write $\req_\val$ when the assignment matters. Equations~\eqref{eq:activation} and~\eqref{eq:requirement} are the semantics of flow workflows. All four tested workflows are flow workflows: under every complete assignment, the product's own walk at the tested commit lists the documents these equations give (Appendix~\ref{app:proofs} says how we checked).

\begin{proposition}[Soundness relative to admitted facts]
\label{prop:sound}
Let the workflow be a flow workflow over two-valued facts. If $\req_\val(d)\in\{\T,\F\}$, then $\req_{\val'}(d)=\req_\val(d)$ for every completion $\val'$ of $\val$.
\end{proposition}

Whatever the open facts turn out to be, a request stays needed and a document marked \skipd stays unneeded. Learning an open fact's value can turn a held document into a request or into \skipd, but never reverses a request or a \skipd. Two-valued evaluation lacks this property. Reading $\U$ as $\F$, the closed-world assumption \cite{reiter1978closed}, drops documents that an open fact may require; reading $\U$ as $\T$ requests them too early.

\begin{proposition}[Strong Kleene evaluation over-holds]
\label{prop:incomplete}
Even in a flow workflow over two-valued facts, strong Kleene evaluation can hold a document that every completion of $\val$ needs.
\end{proposition}

Strong Kleene evaluation forgets that two branches of a fork meet again. The tenancy defect workflow branches on $\mathit{health}$: to ``Escalate health risk'' if it is true, directly to ``Dated defect timeline'' otherwise, and the first leads to the second. With $\mathit{health}=\U$, Equation~\eqref{eq:activation} makes the timeline step $\U$, and every later document whose own condition is not false is held. Under both completions the step is active (Figure~\ref{fig:reconverge}).

\begin{figure}[t]
\centering
\fig{fig_reconverge}{\linewidth}
\Description{Process graph of the tenancy defect workflow evaluated under strong Kleene logic and under exact deferral.}
\caption{The tenancy defect workflow with $\mathit{health}$ unresolved, mould true, the other facts false and no document on file. Strong Kleene evaluation (left) holds every document after the branch whose own condition is not false; the heating report, the inspection and the repair deadline are not needed because their conditions are false. Exact deferral (right) holds only the medical confirmation, on which the two completions disagree. Because nothing is on file, the lease is requested here; in Figure~\ref{fig:hero} the claims attach the lease, so the product lists it as present and does not request it.}
\label{fig:reconverge}
\end{figure}

\subsection{Exact deferral}
\label{sec:exact}

Exact deferral evaluates each document under every completion \cite{vanfraassen1966singular}, as certain answers do for an incomplete database \cite{imielinski1984incomplete}:
\begin{equation}
\exactreq(d) \;=\;
\begin{cases}
\T & \text{if } \req_{\val'}(d)=\T \text{ for every completion } \val' \text{ of } \val,\\
\F & \text{if } \req_{\val'}(d)=\F \text{ for every completion } \val' \text{ of } \val,\\
\U & \text{otherwise.}
\end{cases}
\label{eq:exact}
\end{equation}
The checklist is formed from $\exactreq$ as from $\req$. An open fact becomes a question when changing its value changes a held document or a step.

By definition, exact deferral holds a document exactly when two completions of $\val$ disagree on it, and a document it requests or marks \skipd has that status under every completion; this needs only that a document's need is defined under each complete assignment.

\begin{proposition}[Exact deferral and strong Kleene evaluation]
\label{prop:exact}
In a flow workflow over two-valued facts, exact deferral never holds a document that strong Kleene evaluation requests or marks \skipd.
\end{proposition}

Exact deferral holds what the open facts can change and nothing else. In Figure~\ref{fig:reconverge} it requests the lease, the defect notice and the later documents, holds the medical confirmation and asks whether the mould has affected anyone's health. It costs one evaluation per completion. Above eight open facts or 256 combinations of answers, the product returns the plain walk instead, which stops at the first transition whose condition is open and holds each document whose own condition is open. In Study~3, tenancy claims had at most \sthreeDeferralCasePathTenMaxOpen{} open facts and got exact deferral. Exact deferral never ran on theft: every theft text had \sthreeDeferralCasePathThMinOpenPlain--\sthreeDeferralCasePathThMaxOpen{} open facts and got the plain walk, which on the theft workflow gives what exact deferral gives (Appendix~\ref{app:proofs}). Study~1's research implementation fell back to strong Kleene evaluation above twelve open facts, which its claims never reached.

\subsection{Reading errors}
\label{sec:fate}

A reading can err in two ways: it can leave a settled fact unresolved or give it the wrong value. Let $\val^\ast$ be the true, complete assignment. An \emph{unsafe error} is a request for a document that $\val^\ast$ does not need, or a document marked \skipd that $\val^\ast$ needs.

\begin{proposition}[Reading errors]
\label{prop:errors}
Let the workflow be over two-valued facts, and take whether a document is on file as given. (a) If $\val(f)=\val^\ast(f)$ for every fact with $\val(f)\neq\U$, exact deferral makes no unsafe error. (b) If $\val(f)\neq\val^\ast(f)$ for some fact with $\val(f)\neq\U$, every document whose $\exactreq_\val$ is $\T$ or $\F$ but differs from $\req_{\val^\ast}$ is an unsafe error.
\end{proposition}

An unresolved value costs holds and a question; a wrong value that passes admission produces an error the evaluation cannot detect. A presumption that the file's facts contradict is such a wrong value. Without a model, we computed exactly the expected outcome over the \propClaims{} tenancy claims of Study~1 when each fact is misjudged independently with probability \propEps, once with every misjudgement unresolved and once with every misjudgement the opposite value (Appendix~\ref{app:proofs}, Figure~\ref{fig:fate}). Unresolved values became expected holds under exact deferral and no unsafe error; wrong values became expected missed documents and unnecessary requests.

Admission and agreement exist to turn wrong values into unresolved ones. In Study~3, an analysis added after review found that over \sthreeRThreeCasePathFactRuns{} tenancy fact values (facts $\times$ claim-runs) \sthreeRThreeCasePathAdmittedWrong{} were admitted with a wrong value, each time with all readings agreeing. Of these, \sthreeRThreeCasePathDeniedEventAdmWrong{} admitted a denied defining event as having happened, out of \sthreeRThreeCasePathDeniedEventN{} such claim-runs, in which the readings differed \sthreeRThreeCasePathDeniedEventConflicting{} more times. In the decomposition of Section~\ref{sec:test-cells}, reading errors on the denied-event claims make up about half (\spostDecTenLunaDeniedReadingSharePct\%) of \casepath's strict unsafe tenancy errors (Supplement, Table~\ref{S-tab:fact-accuracy}).
